% GENERATED FILE. Edit canonical lessons, then run npm run generate:books.
% canonical-source-hash: fnv1a64:fef5eea5f1faadeb
% canonical-lessons: SA-C90-ahah, SA-C90-pradosah, SA-C90-aparahnah, SA-C90-madhyaratrih, SA-C90-usah

\chapter{Day, Early evening, Afternoon, Midnight, Dawn}
\label{ch:sa-day-early-evening-afternoon-midnight-dawn}
\hlchaptermodality{90}

\begin{chapteropening}
\textbf{By the end of this chapter:} \emph{I can say a day, early evening, afternoon, midnight and dawn.}

Complete the last lesson of chapter 90: I can say a day, early evening, afternoon, midnight and dawn.
\end{chapteropening}

\section[\texorpdfstring{\sk{अहः} (ahaḥ)}{अहः (ahaḥ)}]{\sk{अहः} (ahaḥ) — a day, daytime}

\label{lesson:SA-C90-ahah}



\begin{quote}
Before the new one: say the Sanskrit for hard, difficult, then the Sanskrit for sweet.
\end{quote}

\subsection*{You'll want to know: \sk{अहः}}
\textbf{\sk{अहः}} — \emph{ahaḥ} — \textquotedblleft{}a day, daytime\textquotedblright{}.

\begin{grammarlens}[title={What you've built}]
One more word for this chapter.
\end{grammarlens}

\subsection*{Your turn}
\begin{itemize}
\raggedright
  \item \emph{Say it:} \emph{ahaḥ}
  \item \emph{Say it:} \emph{ahaḥ}, once more
  \item \emph{Say it:} say \emph{ahaḥ}
  \item \emph{Recall:} say the Sanskrit for hard, difficult, then the Sanskrit for sweet, then say \emph{ahaḥ} again
\end{itemize}

\begin{tcolorbox}[breakable,colback=teal!4,colframe=teal!35!black,title={Before you move on}]
What does \sk{अहः} mean? (\textquotedblleft{}A day, daytime\textquotedblright{}.) Say it once more.
\end{tcolorbox}
\section[\texorpdfstring{\sk{प्रदोषः} (pradoṣaḥ)}{प्रदोषः (pradoṣaḥ)}]{\sk{प्रदोषः} (pradoṣaḥ) — early evening}

\label{lesson:SA-C90-pradosah}



\begin{quote}
Before the new one: say the Sanskrit for sweet, then the Sanskrit for a day.
\end{quote}

\subsection*{You'll want to know: \sk{प्रदोषः}}
\textbf{\sk{प्रदोषः}} — \emph{pradoṣaḥ} — \textquotedblleft{}early evening\textquotedblright{}.

\begin{grammarlens}[title={What you've built}]
One more word for this chapter.
\end{grammarlens}

\subsection*{Your turn}
\begin{itemize}
\raggedright
  \item \emph{Say it:} \emph{pradoṣaḥ}
  \item \emph{Say it:} \emph{pradoṣaḥ}, once more
  \item \emph{Say it:} say \emph{pradoṣaḥ}
  \item \emph{Recall:} say the Sanskrit for sweet, then the Sanskrit for a day, then say \emph{pradoṣaḥ} again
\end{itemize}

\begin{tcolorbox}[breakable,colback=teal!4,colframe=teal!35!black,title={Before you move on}]
What does \sk{प्रदोषः} mean? (\textquotedblleft{}Early evening\textquotedblright{}.) Say it once more.
\end{tcolorbox}
\section[\texorpdfstring{\sk{अपराह्णः} (aparāhṇaḥ)}{अपराह्णः (aparāhṇaḥ)}]{\sk{अपराह्णः} (aparāhṇaḥ) — afternoon}

\label{lesson:SA-C90-aparahnah}



\begin{quote}
Before the new one: say the Sanskrit for a day, then the Sanskrit for early evening.
\end{quote}

\subsection*{You'll want to know: \sk{अपराह्णः}}
\textbf{\sk{अपराह्णः}} — \emph{aparāhṇaḥ} — \textquotedblleft{}afternoon\textquotedblright{}.

\begin{grammarlens}[title={What you've built}]
One more word for this chapter.
\end{grammarlens}

\subsection*{Your turn}
\begin{itemize}
\raggedright
  \item \emph{Say it:} \emph{aparāhṇaḥ}
  \item \emph{Say it:} \emph{aparāhṇaḥ}, once more
  \item \emph{Say it:} say \emph{aparāhṇaḥ}
  \item \emph{Recall:} say the Sanskrit for a day, then the Sanskrit for early evening, then say \emph{aparāhṇaḥ} again
\end{itemize}

\begin{tcolorbox}[breakable,colback=teal!4,colframe=teal!35!black,title={Before you move on}]
What does \sk{अपराह्णः} mean? (\textquotedblleft{}Afternoon\textquotedblright{}.) Say it once more.
\end{tcolorbox}
\section[\texorpdfstring{\sk{मध्यरात्रिः} (madhyarātriḥ)}{मध्यरात्रिः (madhyarātriḥ)}]{\sk{मध्यरात्रिः} (madhyarātriḥ) — midnight}

\label{lesson:SA-C90-madhyaratrih}



\begin{quote}
Before the new one: say the Sanskrit for early evening, then the Sanskrit for afternoon.
\end{quote}

\subsection*{You'll want to know: \sk{मध्यरात्रिः}}
\textbf{\sk{मध्यरात्रिः}} — \emph{madhyarātriḥ} — \textquotedblleft{}midnight\textquotedblright{}.

\begin{grammarlens}[title={What you've built}]
One more word for this chapter.
\end{grammarlens}

\subsection*{Your turn}
\begin{itemize}
\raggedright
  \item \emph{Say it:} \emph{madhyarātriḥ}
  \item \emph{Say it:} \emph{madhyarātriḥ}, once more
  \item \emph{Say it:} say \emph{madhyarātriḥ}
  \item \emph{Recall:} say the Sanskrit for early evening, then the Sanskrit for afternoon, then say \emph{madhyarātriḥ} again
\end{itemize}

\begin{tcolorbox}[breakable,colback=teal!4,colframe=teal!35!black,title={Before you move on}]
What does \sk{मध्यरात्रिः} mean? (\textquotedblleft{}Midnight\textquotedblright{}.) Say it once more.
\end{tcolorbox}
\section[\texorpdfstring{\sk{उषः} (uṣaḥ)}{उषः (uṣaḥ)}]{\sk{उषः} (uṣaḥ) — dawn}

\label{lesson:SA-C90-usah}



\begin{quote}
Before the new one: say the Sanskrit for afternoon, then the Sanskrit for midnight.
\end{quote}

\subsection*{You'll want to know: \sk{उषः}}
\textbf{\sk{उषः}} — \emph{uṣaḥ} — \textquotedblleft{}dawn\textquotedblright{}.

\begin{grammarlens}[title={What you've built}]
One more word for this chapter.
\end{grammarlens}

\subsection*{Your turn}
\begin{itemize}
\raggedright
  \item \emph{Say it:} \emph{uṣaḥ}
  \item \emph{Say it:} \emph{uṣaḥ}, once more
  \item \emph{Say it:} say \emph{uṣaḥ}
  \item \emph{Recall:} say the Sanskrit for afternoon, then the Sanskrit for midnight, then say \emph{uṣaḥ} again
\end{itemize}

\begin{tcolorbox}[breakable,colback=teal!4,colframe=teal!35!black,title={Before you move on}]
What does \sk{उषः} mean? (\textquotedblleft{}Dawn\textquotedblright{}.) Say it once more.
\end{tcolorbox}
